\documentclass[12pt]{article}
\usepackage[utf8]{inputenc}
\usepackage[T1]{fontenc}
\usepackage[margin=1in]{geometry}
\usepackage{setspace}
\usepackage{xcolor}
\usepackage[colorlinks=true, urlcolor=blue!60!black, linkcolor=blue!60!black]{hyperref}
\setlength{\parindent}{2em}
\setlength{\parskip}{1.25em}
\renewcommand{\baselinestretch}{1.0}
\begin{document}
\begin{center}
{\LARGE \textit{SEC v. Jarkesy}, 603 U.S. 109 (2024)\par}
\vspace{0.8em}
{\large Case Brief --- Four-Page Edition\par}
\end{center}
\vspace{1em}

\section{Memory Jogger}

When the SEC seeks civil penalties for securities fraud, the defendant is entitled to a jury trial in an Article~III court: such penalties are a legal remedy for a claim rooted in common-law fraud, and the public rights exception does not let Congress route them to in-house agency adjudication (\href{https://supreme.justia.com/cases/federal/us/603/22-859/}{603 U.S.\ 109}).

\section{Detailed Case Facts}

Between 2007 and 2010, George R.\ Jarkesy, Jr., launched two investment funds that raised about \$24 million from 120 accredited investors, with Patriot28, LLC as adviser. Before 2010 the SEC could seek civil penalties against unregistered entities only in federal court. The Dodd-Frank Act of 2010 let it seek penalties of up to \$725,000 per violation either in court or in its own administrative proceedings. In court, an Article~III judge presides, a jury finds facts, and the Federal Rules of Evidence apply. In-house, there is no jury; the Commission or an ALJ finds facts under rules that permit hearsay, and courts treat those findings as ``conclusive'' if supported by substantial evidence.

\section{Procedural History}

In 2013 the SEC brought an in-house action against Jarkesy and Patriot28 under the antifraud provisions of the 1933, 1934, and 1940 Acts, alleging misrepresentations about investment strategy, the funds' auditor and prime broker, and the funds' value. An ALJ found liability in 2014, and in 2020 the Commission imposed a \$300,000 penalty, a cease-and-desist order, industry bars, and \$685,000 in disgorgement. A divided Fifth Circuit vacated on Seventh Amendment, nondelegation, and Article~II removal grounds. The Supreme Court granted certiorari, heard argument on November 29, 2023, and affirmed on June 27, 2024.

\section{Judicial Votes}

Affirmed and remanded, \textbf{6--3}.
\begin{itemize}
\item \textbf{Majority:} Chief Justice Roberts, joined by Justices Thomas, Alito, Gorsuch, Kavanaugh, and Barrett.
\item \textbf{Concurring:} Justice Gorsuch, joined by Justice Thomas.
\item \textbf{Dissenting:} Justice Sotomayor, joined by Justices Kagan and Jackson.
\end{itemize}

\section{Holding}

When the SEC seeks civil penalties for securities fraud, the Seventh Amendment entitles the defendant to a jury trial in an Article~III court. Under \textit{Tull} and \textit{Granfinanciera}, a statutory claim implicates the Seventh Amendment if it is ``legal in nature,'' and the punitive nature of civil penalties is ``all but dispositive''; Congress may not assign such a claim to an agency unless it falls within a narrow, historically grounded public-rights category. The Court did not reach the nondelegation or removal questions.

\section{Analysis of Opinions}

\textbf{Majority (Roberts).} \emph{Step one:} the remedy is legal because the penalties are designed to punish and deter rather than restore the status quo; they turn on culpability, deterrence, and scienter-based tiers rather than victim loss, and the SEC need not pay them to investors. The cause of action also resembles common-law fraud, and Congress borrowed terms like ``fraud'' and ``deceit'' that bring their ``old soil'' with them. \emph{Step two:} under \textit{Murray's Lessee} and \textit{Stern v.\ Marshall}, Congress may not withdraw from the courts matters that are the subject of a suit at common law. The public rights exception covers historic sovereign prerogatives such as revenue, customs, immigration, tribal relations, public lands, and public benefits, and the government's role as plaintiff does not trigger it. \textit{Atlas Roofing} was distinguished because OSHA's workplace standards were ``unknown to the common law.''

\textbf{Gorsuch (concurring).} The jury right works with Article~III and due process to secure a fair tribunal. He compared SEC in-house proceedings, which the SEC reportedly won about 90\% of the time versus 69\% in court, to colonial vice-admiralty courts and called \textit{Atlas Roofing} an outlier.

\textbf{Sotomayor (dissenting).} Calling the decision a ``power grab,'' she argued that the threshold question is Article~III, that a government civil-penalty claim is the ``heartland'' of public rights under \textit{Murray's Lessee}, \textit{Crowell}, and a unanimous \textit{Atlas Roofing}, and that the ruling imperils more than 200 statutes across two dozen agencies.

\textbf{Later development.} In \textit{FCC v.\ AT\&T} (2026), the Court, 8--1, upheld FCC forfeitures because collection requires a federal suit with a jury trial \textit{de novo}: ``Before a regulated party can be made to pay, the jury gets the last word.''

\section{Examples: Future Applications}

\begin{enumerate}
\item \textbf{Same-Side (Jury Required): Commodities Fraud Penalties.} A punitive in-house penalty against a trader who lied to customers resembles common-law fraud, and commodity regulation is not a historic public-rights category.
\item \textbf{Same-Side (Jury Required): Relabeled Licensing Review.} Calling fraud fines ``licensing reviews'' before a new board fails, because substance controls over labels.
\item \textbf{Opposite-Side (No Jury): Immigration Carrier Fines.} Administrative fines on shipping lines for bringing diseased passengers fall within immigration, a historic public right, as in \textit{Oceanic Steam Navigation Co.\ v.\ Stranahan}.
\item \textbf{Opposite-Side (No Jury): Recovering Benefit Overpayments.} Recouping overpaid benefits involves a public-benefits category and restores the status quo without punishment.
\item \textbf{Fence-Sitter: Immediately Enforceable EPA Penalties.} Punitive emission penalties collectible by offset look like \textit{Jarkesy} and resemble nuisance, but numerical emission thresholds resemble \textit{Atlas Roofing}'s technical standards ``unknown to the common law.''
\end{enumerate}

\section{Critique}

\textbf{Doctrinal tensions.} Critics note that calling the remedy ``all but dispositive'' sits uneasily with preserving \textit{Atlas Roofing}, whose OSHA penalties were equally punitive, leaving that precedent on razor-thin footing. By declining to articulate a comprehensive public-rights test, the Court left lower courts struggling to classify modern regulatory wrongs, such as environmental, labor, or data-privacy violations, that only loosely resemble common-law torts.

\textbf{Defenses.} Proponents praise \textit{Jarkesy}, especially as refined by \textit{FCC v.\ AT\&T}, for restoring a clear baseline: agencies may investigate and issue preliminary assessments, but the Executive cannot unilaterally impose binding punitive fines for common-law-like wrongs without an Article~III jury.

\section{Key Quotations}

\begin{quote}
``In this case, the remedy is all but dispositive. \ldots\ In sum, the civil penalties in this case are designed to punish and deter, not to compensate. They are therefore `a type of remedy at common law that could only be enforced in courts of law.' ''
\end{quote}
\noindent --- Chief Justice Roberts, Majority Opinion, 603 U.S.\ at 123, 125.

\begin{quote}
``What matters is the substance of the suit, not where it is brought, who brings it, or how it is labeled.''
\end{quote}
\noindent --- Chief Justice Roberts, Majority Opinion, 603 U.S.\ at 135.

\begin{quote}
``A defendant facing a fraud suit has the right to be tried by a jury of his peers before a neutral adjudicator. Rather than recognize that right, the dissent would permit Congress to concentrate the roles of prosecutor, judge, and jury in the hands of the Executive Branch. That is the very opposite of the separation of powers that the Constitution demands.''
\end{quote}
\noindent --- Chief Justice Roberts, Majority Opinion, 603 U.S.\ at 140--141.

\end{document}
